% academic-article.tex：标准学术文档骨架模板
% 编译命令：latexmk -pdf main.tex（无 bib 时编译两遍即可）
% 用途：任务未指定模板时的默认骨架。指定模板时以模板为准。
\documentclass[11pt,a4paper]{article}

% 宏包按安全顺序排列：先文字编码，再数学，再图表，最后工具
\usepackage[utf8]{inputenc}      % pdflatex 必需；xelatex 下删除本行
\usepackage[T1]{fontenc}         % pdflatex 必需；xelatex 下删除本行
\usepackage{amsmath,amssymb}     % 数学
\usepackage{graphicx}            % 插图
\usepackage{booktabs}            % 三线表
\usepackage{multirow}            % 合并单元格
\usepackage{tabularx}            % 自适应列宽，防溢出
\usepackage{hyperref}            % 最后加载超链接

\title{论文标题}
\author{作者姓名}
\date{\today}

\begin{document}

\maketitle

\begin{abstract}
摘要段落。只写材料中出现的事实，不新增材料外内容。
\end{abstract}

\section{引言}\label{sec:intro}
正文。引用规范写法：见文献~\cite{key2024}。交叉引用规范写法：见第~\ref{sec:method}节。

\section{方法}\label{sec:method}
\subsection{模型定义}
\begin{equation}\label{eq:model}
  f(x) = \sum_{i=1}^{n} w_i \, \phi(x_i)
\end{equation}
公式~\eqref{eq:model} 定义了模型。多行公式用 align：
\begin{align}
  a &= b + c \label{eq:first}\\
  d &= e - f \label{eq:second}
\end{align}
变量用斜体（默认），单位与算子用正体：$v = 5\,\mathrm{m/s}$。

\subsection{定理环境示例}
\newtheorem{theorem}{定理}
\begin{theorem}\label{thm:main}
定理陈述。
\end{theorem}
\begin{proof}
证明内容。
\end{proof}

\section{实验}\label{sec:exp}
图规范写法（必须有 caption 与 label，正文用 \ref 引用）：
\begin{figure}[htbp]
  \centering
  \includegraphics[width=0.8\linewidth]{fig1}
  \caption{图标题。}\label{fig:demo}
\end{figure}

表规范写法（booktabs 三线表，禁止竖线）：
\begin{table}[htbp]
  \centering
  \caption{表标题。}\label{tab:demo}
  \begin{tabular}{lcc}
    \toprule
    方法 & 指标A & 指标B \\
    \midrule
    基线 & 0.85 & 0.72 \\
    本文 & 0.91 & 0.79 \\
    \bottomrule
  \end{tabular}
\end{table}

宽表用 tabularx 防溢出：
\begin{table}[htbp]
  \centering
  \begin{tabularx}{\linewidth}{lX}
    \toprule
    项目 & 说明（自动换行列） \\
    \midrule
    条目 & 较长的说明文字放在 X 列。 \\
    \bottomrule
  \end{tabularx}
\end{table}

\section{结论}\label{sec:conclusion}
结论段落。

% 参考文献二选一：
% 方式一（内嵌 bibliography，无 bib 文件时使用）：
\begin{thebibliography}{9}
\bibitem{key2024} 作者. 标题. 期刊, 2024.
\end{thebibliography}

% 方式二（bib 文件，注释方式一并改为 \bibliography{refs}）：
% \bibliographystyle{plain}
% \bibliography{refs}

\end{document}
